% Resumo em Português
\setlength{\absparsep}{18pt}

\begin{resumo}

<<<<<<< HEAD
A adoção do Building Information Modeling (BIM) vem transformando digitalmente a construção civil ao centralizar informações de todo o ciclo de vida de um projeto, aumentando a precisão, reduzindo erros e otimizando recursos. Esse crescimento, contudo, traz desafios relacionados ao gerenciamento de grandes volumes de dados e, sobretudo, à integração de normas regulatórias complexas. No setor de Arquitetura, Engenharia e Construção (AEC), a verificação automática de conformidade ainda esbarra na linguagem técnica e na falta de padronização dos documentos normativos. Uma alternativa promissora é a metodologia RASE (\textit{Requirement, Applicability, Selection, Exception}), que converte normas em dados computáveis. Em paralelo, Modelos de Linguagem de Grande Porte (LLMs) têm sido investigados para interpretar textos regulatórios, mas ainda há escassez de avaliações sistemáticas sobre LLMs \textit{open-source} aplicados a normas brasileiras. Esta pesquisa avalia comparativamente seis LLMs \textit{open-source} --- Llama, Dolphin, Gemma, Mistral, Alpaca e Qwen --- na conversão automática de normas técnicas da AEC para o formato RASE, utilizando \textbf{exclusivamente Engenharia de Prompt} (sem \textit{Fine-Tuning} e sem \textit{Retrieval-Augmented Generation}). A tarefa é decomposta em três níveis sucessivos de normalização: N1 --- segmentação atômica das sentenças normativas; N2 --- identificação dos operadores RASE (Requisito, Aplicabilidade, Seleção, Exceção); e N3 --- estruturação semântica em JSON. Foram conduzidos três experimentos: EN1 (N1 isolado), EN2 (N2 isolado, a partir do N1 de referência) e EN1N2 (pipeline encadeado), permitindo analisar a propagação de erros entre as etapas. Os modelos foram servidos localmente via Ollama e aplicados a um dataset com 79 normas brasileiras anotadas, baseado na NBR 9050. A qualidade das saídas foi avaliada por treze métricas, organizadas em cinco famílias: similaridade lexical (FuzzyWuzzy, TF-IDF), semântica (SBERT, BERTimbau, Multilingual), de distância (\textit{Word Mover's Distance} com FastText e NILC), voltadas para geração (BERTScore e ROUGE-L) e de classificação (Acurácia, Precisão, Revocação e F1-score). Considerando a média global sobre as nove métricas de similaridade nos três experimentos, o Dolphin lidera (0,676), seguido pelo Llama (0,658) e pelo Mistral (0,654); o Dolphin também oferece o melhor custo-benefício, com tempo de execução cerca de noventa vezes menor que o do Llama no pipeline encadeado. O Llama lidera o experimento EN2 e o nível N3 isolado nas métricas semânticas. As métricas semânticas e o BERTScore mostraram-se mais adequadas que as lexicais para avaliar transformações normativas, com o BERTimbau apresentando ganho consistente sobre o SBERT em português brasileiro. As métricas de classificação revelaram que os operadores Seleção e Exceção são as classes mais difíceis em N2, com F1 próxima de zero para todos os modelos. A pesquisa contribui com a definição operacional dos três níveis de normalização sobre a metodologia RASE, com um conjunto reproduzível de prompts especializados por nível e operador, e com a primeira comparação sistemática de seis LLMs \textit{open-source} para a conversão de normas brasileiras da AEC.
=======
A verificação automática de conformidade de normas técnicas no setor de Arquitetura, Engenharia e Construção (AEC) ainda é limitada pela linguagem técnica e pela falta de padronização dos documentos normativos. A metodologia RASE (\textit{Requirement, Applicability, Selection, Exception}) converte normas em dados computáveis, e Modelos de Linguagem de Grande Porte (LLMs) têm sido investigados para apoiar essa conversão. Esta pesquisa avalia comparativamente seis LLMs \textit{open-source} --- Llama, Dolphin, Gemma, Mistral, Alpaca e Qwen --- na conversão automática de normas brasileiras da AEC para o formato RASE, utilizando \textbf{exclusivamente Engenharia de Prompt}. A tarefa é decomposta em três níveis (N1: segmentação atômica; N2: identificação dos operadores RASE; N3: estruturação semântica em JSON) e avaliada em cinco experimentos (EN1, EN2, EN1N2, EN3 e EN1N2N3) sobre um \textit{dataset} de 79 normas anotadas baseado na NBR 9050. A qualidade foi medida por métricas de similaridade lexical, semântica e de distância, complementadas por métricas de classificação. O Llama lidera a qualidade global, com similaridade semântica superior a 0,80 em EN2, enquanto o Dolphin apresenta o melhor custo-benefício ao combinar qualidade competitiva com tempo de execução até cem vezes menor. As métricas semânticas mostraram-se mais adequadas que as lexicais, e o pipeline encadeado evidenciou a propagação esperada de erros entre os níveis.
>>>>>>> b1bd44c (citation)

\textbf{Palavras-chave}: Aprendizado de Máquina, Processamento de Linguagem Natural, LLM, Engenharia de Prompt, BIM, RASE.

\end{resumo}
